\section{Elementary-function certificates that preserve exact zeros}
\label{sec:analytic}
\status{Implemented in Python for one real variable: rational interval evaluation of a fixed elementary-expression grammar, checked coverage trees, exact jets, and bounded automatic anchor/order discovery. General Taylor-model arithmetic, multivariate certificates, and Lean soundness proofs are proposed extensions.}

\subsection{Two obstacles, not one}
A polynomial certificate engine can reason about \(u=\exp(x)\) as a fresh variable, but it cannot obtain useful relationships between \(u\) and \(x\) merely by expanding polynomials. A semantic bridge must supply proved elementary-function facts.

A second obstacle remains even after adding rigorous numerical bounds. Suppose the goal is
\[
 f(x)\ge0\quad(x\in B),
\]
and \(f\) actually vanishes at an interior point. An enclosure of the form \([f(x)-\varepsilon,f(x)+\varepsilon]\), however small its positive error, does not prove nonnegativity at that zero. Subdividing around the point may leave a leaf containing it forever. The appropriate remedy is to expose an exact vanishing factor or derivative structure, not to interpret a sufficiently small negative lower bound as zero.

Verified Taylor/interval methods are established in other proof assistants. Solovyev and Hales describe formal verification of nonlinear inequalities with Taylor interval approximations, while the Rocq Interval package supplies an existing proof-producing interval methodology~\cite{solovyev-hales,rocq-interval}. The proposal here is a Forge adapter with a deliberately small initial rational grammar and a zero-aware certificate constructor. It is not a claim to invent validated numerics.

\subsection{The admitted grammar}
The executable accepts expression trees generated by
\[
 e::=c\mid x\mid -e\mid e+e\mid ee\mid e/e\mid e^n
       \mid\exp(e)\mid\log(e)\mid\sin(e)\mid\cos(e),
\]
where \(c\in\Q\) and \(n\in\N\) is a literal. The problem is an original expression and a closed rational interval \(B=[\ell,u]\), asking for \(e(x)\ge0\) throughout \(B\).

Every division requires an interval denominator excluding zero. Every logarithm requires a strictly positive interval argument. These are admission conditions for the analytic lemmas, not assertions that unsupported source expressions are false. In particular, Lean's use of total functions does not remove the side conditions of the real-analytic theorems used to justify differentiation and series bounds.

The checker traverses the \emph{original} expression tree on every leaf, even for a jet proof whose derivative has simplified dramatically. This prevents a simplifier from erasing the domain of a subexpression. The explicit expression \(0/x\) on \([-1,1]\) is rejected by this analytic fragment: rewriting a displayed numerator to zero does not establish the nonvanishing premise of the reciprocal derivative rules.

The implementation is univariate. Expressions may be nested, but there is only one free real coordinate. A multivariate natural-interval extension is straightforward at the representation level; the higher-order derivative and vanishing certificates require additional care and are not claimed to exist in the supplied code.

\subsection{Exact rational interval arithmetic}
For intervals \(A=[a_-,a_+]\) and \(B=[b_-,b_+]\), the checker implements
\begin{align*}
 A+B&=[a_-+b_-,a_++b_+],\\
 -A&=[-a_+,-a_-],\\
 AB&=[\min S,\max S],
 \quad S=\{a_-b_-,a_-b_+,a_+b_-,a_+b_+\}.
\end{align*}
If \(0\notin B\), then
\[
 B^{-1}=[1/b_+,1/b_-].
\]
Integer powers use endpoint monotonicity; an even positive power has lower bound zero when its argument interval crosses zero. The zeroth power is exactly \([1,1]\).

All endpoints are rational numbers represented by arbitrary-precision integer numerator and positive denominator. There is no floating-point rounding mode to configure in the checker. This makes exact replay simple, though potentially expensive when high-order expansions create large denominators. Bit-length and operation budgets remain necessary in a production implementation.

\subsection{A small library of rational elementary bounds}
The first implementation uses expansions around zero, except for a transformed logarithm. They are intentionally elementary enough to prove independently and to recompute without calling a transcendental numerical library.

\subsubsection{Exponential}
For a rational point \(a\), let
\[
 T_N(a)=\sum_{k=0}^N\frac{a^k}{k!},\qquad r=|a|.
\]
If \(r<N+2\), define
\begin{equation}\label{eq:exp-error}
 E_N(r)=\frac{r^{N+1}}{(N+1)!}
         \frac{1}{1-r/(N+2)}.
\end{equation}
Then
\begin{equation}\label{eq:exp-enclosure}
 T_N(a)-E_N(r)\le\exp(a)\le T_N(a)+E_N(r).
\end{equation}
\begin{proof}
The absolute value of the omitted term of index \(N+1\) is \(r^{N+1}/(N+1)!\). Every subsequent ratio is at most \(r/(N+2)<1\). The absolute tail is therefore bounded by the geometric series in \eqref{eq:exp-error}. Applying this to the exponential power series proves \eqref{eq:exp-enclosure}.
\end{proof}
For an interval argument \([a,b]\), monotonicity gives a lower bound from the point enclosure at \(a\) and an upper bound from the point enclosure at \(b\). The producer computes the Taylor polynomial by powers and factorials. The checker recomputes it using a term recurrence, retaining the same mathematically justified tail bound.

\subsubsection{Logarithm}
For \(a>0\), set
\[
 t=\frac{a-1}{a+1},\qquad |t|<1.
\]
Use
\begin{equation}\label{eq:log-series}
 \log a=2\sum_{k=0}^{M-1}\frac{t^{2k+1}}{2k+1}+R_M,
 \qquad
 |R_M|\le
 \frac{2|t|^{2M+1}}{(2M+1)(1-t^2)}.
\end{equation}
\begin{proof}
Integrate the geometric expansion of \(1/(1-s^2)\), or equivalently use the odd power series for \(2\operatorname{artanh}(t)\). In the absolute tail, each denominator is at least \(2M+1\); bounding it by that denominator leaves a geometric sum in \(|t|^2\). The identity \(a=(1+t)/(1-t)\) gives the logarithm formula.
\end{proof}
Monotonicity again reduces an interval argument to two rational endpoint computations. No rational approximation to a logarithmic constant is trusted as an input; it is reconstructed by this rule.

\subsubsection{Sine and cosine}
Take the Maclaurin polynomial of degree \(N\), with zero coefficients of the wrong parity. On an argument interval \(A\), let \(r=\max_{a\in A}|a|\). Since every derivative of sine or cosine has absolute value at most one,
\begin{equation}\label{eq:trig-error}
 |f(a)-T_N(a)|\le\frac{r^{N+1}}{(N+1)!}
 \qquad(a\in A).
\end{equation}
The polynomial is evaluated by exact interval arithmetic and enlarged by that rational error. Intersecting with the known range \([-1,1]\) is sound and often useful.

The last intersection explains a useful automatic proof improvement. To prove \(1-\cos x\ge0\), a checker does not need the Taylor polynomial to resolve the exact zero at \(x=0\): the universal upper bound \(\cos x\le1\) already preserves it. The jet selector can exploit this fact when a higher derivative has that form.

\subsection{Coverage certificates rather than sampled grids}
A certificate tree starts from the problem's original interval. A leaf records a series order and an advertised rational enclosure. The checker recomputes an enclosure of the original expression on that leaf, and accepts only if the advertised enclosure contains the recomputed one and has nonnegative lower endpoint.

An internal node records a rational split point \(s\) strictly between the current endpoints. Its children are interpreted as \([\ell,s]\) and \([s,u]\). The checker derives these child domains itself. It does not trust independently supplied child boxes, so there is no opportunity to omit a small interval around a troublesome point.

\begin{theorem}[Interval-tree soundness]\label{thm:interval}
Assume the elementary atom bounds are valid and the interval operations preserve enclosure. An accepted coverage tree proves \(e(x)\ge0\) for every \(x\) in the original interval, with all required analytic-domain conditions satisfied there.
\end{theorem}
\begin{proof}
Structural induction on the expression gives enclosure on each checked leaf. Its nonnegative lower endpoint proves the local inequality. Induction on the tree combines the two child results, using the exact interval-cover identity at every split. The original-expression domain traversal supplies the side conditions on all leaves and hence on their union.
\end{proof}

The default proposer tries orders \(8,16,32\), then bisects the interval if needed. It has a depth and node budget. A finite precision ceiling is not a completeness guarantee. A certificate produced by another search policy can still be accepted, provided its recorded orders and tree fit the checker's explicit limits.

The workload includes twelve positive-margin expressions of the form \(x-x+\varepsilon\). The input AST deliberately retains both occurrences of \(x\), so natural interval arithmetic gives \([\ell-u+\varepsilon,u-\ell+\varepsilon]\). Subdivision is needed until the leaf width is small enough. These are transparent coverage tests, not hard mathematical inequalities; a symbolic simplifier would close them immediately.

\subsection{Vanishing jets: remove the uncertain value at the zero}
A jet certificate records an anchor \(a\), an order \(m\ge1\), exact vanishing conditions, and an interval-tree proof of a sign condition on the \(m\)-th derivative.

\begin{theorem}[Zero-aware nonnegativity]\label{thm:jet}
Let \(f\) be \(C^m\) on a neighborhood of a closed interval \(B\), let \(a\in B\), and suppose
\begin{equation}\label{eq:vanishing}
 f(a)=f'(a)=\cdots=f^{(m-1)}(a)=0,
 \qquad f^{(m)}(x)\ge0\quad(x\in B).
\end{equation}
If \(m\) is even, then \(f(x)\ge0\) throughout \(B\). If \(m\) is odd, the same conclusion holds on \(B\cap[a,\infty)\).
\end{theorem}
\begin{proof}
At \(x=a\), the conclusion is an exact vanishing condition. Otherwise Taylor's theorem gives a point \(\xi\) strictly between \(a\) and \(x\) such that
\[
 f(x)=\frac{f^{(m)}(\xi)}{m!}(x-a)^m.
\]
The derivative factor is nonnegative and the factorial is positive. The power is nonnegative for even \(m\), or for \(x\ge a\) with arbitrary \(m\). Multiplying gives the claim.
\end{proof}

The prototype's odd-order rule requires the whole input interval to lie to the right of the anchor. A left-sided odd-order certificate could instead use \(f^{(m)}\le0\); that symmetric constructor is not implemented. Rejecting an odd right-oriented certificate on \([-1,1]\) for \(f(x)=x^3\) is one of the unit tests.

The exact-zero condition is stronger than an interval containing zero. The anchor evaluator accepts rational arithmetic together with the exact identities
\[
 \exp(0)=1,\quad\log(1)=0,\quad\sin(0)=0,\quad\cos(0)=1.
\]
If an anchor expression cannot be evaluated by those rules, the attempt is inconclusive. A high-precision approximation close to zero is never converted to an exact jet coefficient.

\subsection{Four worked certificates}

\paragraph{Exponential tangent.}
For \(f(x)=e^x-1-x\), the anchor values \(f(0)=f'(0)=0\) are exact and \(f''(x)=e^x\). On \([-1,1]\), a rational exponential lower enclosure is positive. The order-two certificate proves
\[
 e^x\ge1+x.
\]
The point \(x=0\) is handled by equality, not by a numerical tolerance.

\paragraph{Logarithmic tangent.}
For \(f(x)=x-\log(1+x)\), the same two anchor coefficients vanish and
\[
 f''(x)=\frac{1}{(1+x)^2}>0
\]
on \([-1/2,1]\). The original domain check proves \(1+x\ge1/2\); the reciprocal interval therefore excludes a pole. This gives
\[
 \log(1+x)\le x.
\]

\paragraph{Cosine lower approximation.}
Set \(f(x)=\cos x-1+x^2/2\). Its first four jet coefficients at zero vanish, and \(f^{(4)}(x)=\cos x\). A positive cosine enclosure on \([-1,1]\) gives the recorded order-four certificate for
\[
 \cos x\ge1-\frac{x^2}{2}.
\]
A shorter proof also exists: \(f''=1-\cos x\ge0\), so order two suffices. The automatic selector finds that lower-order certificate by exploiting the universal trigonometric range bound.

\paragraph{Sine lower approximation.}
Let \(f(x)=\sin x-x+x^3/6\). The first five jet coefficients vanish at zero and \(f^{(5)}(x)=\cos x\). On \([0,1]\), the right-sided order-five rule proves
\[
 \sin x\ge x-\frac{x^3}{6}.
\]
Again a lower-order proof is possible: \(f'''=1-\cos x\ge0\). The automatic selector finds an order-three certificate on the right-sided interval.

The manual-order and automatic-order certificates concern the same four mathematical goals and are not counted as additional independent examples. Their difference is useful nonetheless: it demonstrates that a certificate's jet order need not equal the function's exact order of vanishing. Any lower order satisfying the theorem's hypotheses is sufficient.

\subsection{Automatic proposal without guessing exact zeros}
The additional \code{auto\_prove} entry point first attempts a shallow interval tree. If that fails, it tries the anchor set consisting of zero and the two rational interval endpoints, removing duplicates and anchors outside the domain. At each anchor it differentiates repeatedly, checking exact vanishing after each step, up to a configured maximum order.

If a required coefficient is nonzero or unsupported by the exact evaluator, larger orders at that anchor cannot justify the same vanishing certificate and are not explored. An odd order with an interval extending left of the anchor is skipped. Every admissible order then gets a bounded derivative-sign proof attempt.

\par\noindent\begin{minipage}{\linewidth}
\begin{lstlisting}
try_exact_jets(expression, interval):
  try bounded ordinary interval proof
  for anchor in {0, left endpoint, right endpoint}:
      current := expression
      for order from 1 to 6:
          if exact_value(current, anchor) is not exactly zero:
              stop this anchor
          current := symbolic_derivative(current)
          if orientation is admissible:
              try a bounded interval proof current >= 0
              on success return the jet certificate
  return UNKNOWN
\end{lstlisting}
\end{minipage}\par

This selector was executed on the four worked goals and independently checked. It found orders \(2,2,2,3\), respectively. The count of attempted proof profiles was \(2,2,2,4\), including the failed initial interval attempt. These are recorded separately in \path{results/jet-discovery.json}.

More ambitious anchor proposal can use approximate minimizers or roots of a polynomial surrogate, but approximate zeros must then be turned into exact anchor identities or into isolating-interval arguments. A candidate numerical minimum with value near zero is useful guidance and nothing more. Algebraic anchors would require an exact algebraic-number representation and a bridge to the existing polynomial root machinery.

\subsection{The exact meaning of checker independence}
The search and checker are separate modules. The search uses simplified symbolic derivatives; the checker derives the expression independently with unsimplified differentiation rules. Search-side Taylor sums use powers and factorials; checker-side coefficients are recomputed by recurrences. Both use Python's exact rational arithmetic and the same mathematical atom-bound specifications.

Consequently this is meaningful producer/checker separation, but not a formal correctness proof. A mistaken mathematical bound shared by both implementations could escape these tests. Nor is the checker's unsimplified derivative always as precise as the searcher's simplified interval expression. It may therefore reject a sound candidate whose advertised enclosure is too narrow for its own evaluation. That is a false negative, not permission to widen the accepted semantics. A production reflected derivative DAG with proved simplification would remove much of this mismatch.

The optional numerical crosscheck evaluates 320 rational point arguments with 80 decimal digits using \code{mpmath}. It is explicitly outside acceptance, uses a small tolerance only for comparison at exact numerical endpoints, and is reported as a sanity check rather than mathematical proof. The certificates themselves use no \code{mpmath}, floating-point transcendental routine, or external numeric library.

\subsection{A limited completeness theorem}
\begin{theorem}[Robust strict positivity, idealized search]\label{thm:robust}
Consider an admitted elementary AST on a compact rational interval, with every logarithm argument bounded away from zero and every divisor bounded away from zero. If its denotation is strictly positive throughout the interval, an unbounded fair search combining arbitrarily fine subdivision and arbitrarily high series orders eventually finds a finite interval certificate.
\end{theorem}
\begin{proof}
Continuity and compactness give a positive minimum \(\delta\). The domain margins persist in sufficiently small neighborhoods. At each point, the elementary rational enclosures converge as the order increases, and the natural interval extension converges to the point value as the interval shrinks. By structural induction on the finite AST, a sufficiently small rational neighborhood and sufficiently high order give a lower enclosure above zero. Compactness supplies a finite subcover. A sufficiently fine subdivision refines these neighborhoods, and sufficiently large per-leaf orders establish positivity on the resulting finite partition. Fair exploration eventually reaches such choices.
\end{proof}

This theorem does not apply to arbitrary nonnegative functions with zeros. The jet rule covers a useful additional class but is not a general nonnegativity decision procedure. It also does not apply to the literal fixed order set \(\{8,16,32\}\), fixed depth caps, or fixed jet order limit in the supplied executable. Distinguishing a limiting algorithm from its budgeted implementation is necessary for an honest capability statement.

\subsection{A stronger production lane: polynomial Taylor models}
Natural interval arithmetic loses dependencies between repeated subexpressions. The next step should reuse Forge's existing polynomial engines rather than implement another polynomial positivity checker. On a rational box, represent an expression by
\[
 e(x)=P(x)+\rho(x),\qquad \rho(x)\in I,
\]
where \(P\) is a rational polynomial and \(I\) a rational interval. Addition combines the polynomials and remainder intervals. Multiplication keeps the polynomial product, truncates it only with a separately bounded tail, and bounds the terms involving remainders. Elementary composition expands around a chosen rational center and attaches a proved tail bound.

A goal \(e\ge0\) becomes the polynomial obligation \(P+\inf I\ge0\), which the existing cone, Bernstein, or Sturm workers may settle. Near a zero, a jet or exact factorization should be applied before discarding correlation into an interval remainder. Otherwise a sophisticated Taylor model can reproduce the same zero obstruction as a simple interval evaluator.

This Taylor-model implementation is not included. The current atom enclosures and coverage checker establish a smaller integration target first. Mathlib's real-variable Taylor library already exposes Lagrange and integral remainder results, including \path{taylor_mean_remainder_lagrange_iteratedDeriv}; those are concrete theorem-bridge targets, not a claim that their precise signatures were compiled against the repository's pinned Mathlib revision~\cite{mathlib-taylor}.

\subsection{Library choices and range reduction}
For search, \code{python-flint}'s \code{arb} type is a useful optional interval proposal engine: its documented representation uses a floating-point midpoint and an outward radius, and its precision is configurable~\cite{arb}. A Forge adapter can use such balls to select subdivisions, expansion orders, and centers. It must still export rational bounds or an explicit analytic certificate accepted by the Lean-side checker. An Arb interval is not by itself a Lean proof.

The reference formulas work best for moderate arguments. Large exponential arguments should be reduced by identities such as \(\exp(x)=\exp(x/2^k)^{2^k}\), with repeated squaring checked explicitly. A positive logarithm argument can be scaled into a compact range and accompanied by a separately certified \(\log 2\) term. Trigonometric reduction requires certified bounds involving \(\pi\) and proof that the chosen reduction index is valid; a floating-point quadrant computation alone is not enough. None of these range-reduction optimizations is present in the prototype, so large arguments can legitimately produce \unknown{} under its fixed caps.
